\chapter{Выбор программных и аппаратных средств для реализации приложения}

Для того чтобы создать Android-приложение, способное обнаруживать дефекты 3D-печати, необходимо комплексное решение, объединяющее знания в области материаловедения, машинного зрения и мобильной разработки. Прежде всего, следует провести анализ и классификацию наиболее распространённых дефектов, таких как недостаточная адгезия первого слоя, деформация, перегрев, недостаточное охлаждение, а также «спагетти» \cite{заикин2023проектирование}. Именно эти проблемы чаще всего возникают у новичков, и их распознавание может стать отправной точкой для дальнейшей оптимизации процесса печати.

Затем следует создать алгоритмы компьютерного зрения, которые будут точно распознавать дефекты на основе изображений, получаемых с камеры Android-устройства. Для этого можно применить современные методы обработки изображений и сверточные нейронные сети, специально настроенные на задачу сегментации и классификации дефектов \cite{измайлов2020анализ}.

Следует уделить внимание созданию интуитивно понятного пользовательского интерфейса, который позволит оперативно получать обратную связь о качестве печати и своевременно корректировать настройки 3D-принтера.

\section{Выбор средств для аннотирования изображений и обучения нейросети}

Так как задача создания приложения для обнаружения дефектов требует оперативной и точной идентификации отклонений в режиме реального времени, оптимальным выбором является использование архитектуры YOLO (You Only Look Once, что в переводе с английского означает «ты смотришь только один раз»). YOLO -- это сверточная нейронная сеть, которая обрабатывает изображение за один проход, одновременно определяя местоположение и тип объектов. Такой подход обеспечивает высокую скорость работы и точность детекции, что особенно важно для мобильных решений \cite{брехт2022применение}.

Например, в отличие от двухступенчатых методов, таких как Faster R-CNN (Faster Region-based Convolutional Neural Network -- Быстрая Региональная Сверточная Нейронная Сеть), где сначала генерируются регионы интереса, а затем происходит их классификация, YOLO обрабатывает изображение целиком за один проход. Такой подход позволяет достигать скоростей до 45–155 кадров в секунду, что особенно важно для систем, работающих в режиме реального времени (real-time systems), например, для мобильных приложений по обнаружению дефектов или систем видеонаблюдения.

Еще одним примером является сравнение с архитектурой SSD (Single Shot MultiBox Detector, «Однокадровый мультибокс-детектор»). Несмотря на то, что SSD также реализует детекцию объектов за один проход, YOLO часто демонстрирует более высокую производительность на сложных изображениях благодаря более глубокому и оптимизированному сверточному базису и улучшенной стратегии прогнозирования ограничивающих рамок \cite{lishchenko2022online}.

Кроме того, благодаря своей архитектуре YOLO может эффективно интегрироваться в облачные и мобильные приложения, что делает её универсальным инструментом для широкого спектра сценариев использования.

Для аннотирования изображений был выбран Roboflow -- это облачная платформа для управления наборами данных, разметки и обучения моделей компьютерного зрения. Она предоставляет удобный RESTful API (интерфейс программирования приложений, построенный на принципах архитектуры REST, когда запросы отправляются по протоколу HTTP) и библиотеку для Python. Это упрощает интеграцию процесса подготовки данных и обучения нейросетей в различные приложения.

Основное преимущество Roboflow заключается в том, что она позволяет легко загружать, аннотировать и преобразовывать изображения, а также применять аугментацию -- метод искусственного увеличения объёма обучающего набора данных путём внесения случайных изменений, таких как повороты, масштабирование и изменение яркости \cite{kukartsev2024deep}. Такой подход значительно ускоряет подготовку набора данных для обучения моделей, обеспечивая более высокое качество финальных результатов.

Roboflow, благодаря своей облачной инфраструктуре, может работать с большими объёмами данных, что крайне важно для обучения сложных моделей. Это позволяет выполнять цикл обучения быстро и оптимизировать параметры модели.

После обучения модель можно развернуть и использовать через API. Это даёт возможность быстро внедрять системы обнаружения объектов в мобильные приложения или промышленные системы мониторинга\cite{богданов2019технология}.

\section{Выбор интегрируемой среды разработки для создания мобильного приложения на Android}

При разработке мобильных приложений для Android выбор оптимальной интегрированной среды разработки (Integrated Development Environment, IDE) -- это один из самых важных этапов. От этого выбора зависит, насколько эффективно будет возможно писать, отлаживать и тестировать код.

После тщательного анализа доступных инструментов, таких как Eclipse с плагином Android Development Tools (ADT -- инструменты разработки для Android) и IntelliJ IDEA, особое внимание было уделено Android Studio\cite{kocakoyun2017developing}. Эта среда разработки, являющаяся официальным продуктом компании Google, обеспечивает своевременное обновление и полную совместимость с последними версиями Android Software Development Kit (SDK — набор средств разработки программного обеспечения), что гарантирует доступ к современным функциям и возможностям, необходимым для создания качественных приложений.

Одним из главных преимуществ Android Studio является интеграция с системой сборки Gradle -- мощным инструментом автоматизации, который помогает управлять зависимостями проекта и оптимизировать процесс сборки. Эта интеграция значительно ускоряет цикл разработки, а также способствует повышению стабильности и производительности конечного продукта\cite{craig2015learn}.

После изучения различных инструментов был сделан выбор в пользу Android Studio\cite{zapata2016android}. Поскольку и Android, и Android Studio являются официальными продуктами компании Google, эта среда разработки регулярно обновляется и полностью совместима с последними версиями операционной системы Android.